% Generated draft: requires mathematical and source review.
\documentclass[11pt]{article}
\usepackage[margin=25mm]{geometry}
\usepackage{fontspec}
\setmainfont{Latin Modern Roman}
\usepackage{amsmath,amssymb,mathtools,unicode-math}
\setmathfont{Latin Modern Math}
\usepackage{enumitem,needspace,xurl,hyperref}
\setlength{\parindent}{0pt}
\setlength{\parskip}{5pt}
\setlength{\emergencystretch}{3em}
\newcommand{\N}{\mathbb{N}}
\newcommand{\Z}{\mathbb{Z}}
\newcommand{\Q}{\mathbb{Q}}
\newcommand{\R}{\mathbb{R}}
\newcommand{\C}{\mathbb{C}}
\newcommand{\F}{\mathbb{F}}
\newcommand{\A}{\mathbb{A}}
\newcommand{\PP}{\mathbb{P}}
\newcommand{\E}{\mathbb{E}}
\newcommand{\eps}{\varepsilon}
\newcommand{\dd}{\,\mathrm{d}}
\begin{document}
\section*{UnsolvedMath ID: 30002387}
\textbf{Source problem code:} catalogue entry
\subsection{Definitions and mathematical statement}
Let \(M\) be a compact Kähler manifold of complex dimension \(n\), and identify the Kähler metrics \(\omega\) and \(\chi\) with their smooth, real, closed, positive \((1,1)\)-forms. Write \([\chi]\) for the real de Rham cohomology class of \(\chi\), and write \(\alpha^j\) for the \(j\)-fold wedge product of a form \(\alpha\), with \(\alpha^0=1\).

The supplied record does not specify the range of the integer \(k\). The displayed equation requires \(0\leq k\leq n\); the intended range within these bounds remains to be confirmed. For any such fixed \(k\), integration of the proposed equation forces the normalization
\[
 c_k=\binom{n}{k}\frac{\displaystyle\int_M\chi^{n-k}\wedge\omega^k}{\displaystyle\int_M\chi^n}.
\]
This expression depends only on the cohomology classes of \(\chi\) and \(\omega\). It gives the normalization required by the displayed equation; the supplied record calls \(c_k\) the associated topological constant without explicitly defining it.

\textbf{Existence question.} Given \((M,\omega)\), \(\chi\), and an intended index \(k\) as above, does there exist a smooth Kähler form \(\widetilde\chi\) with \([\widetilde\chi]=[\chi]\) such that the following equality of differential forms holds everywhere on \(M\)?
\[
 c_k\widetilde\chi^n
 =\binom{n}{k}\widetilde\chi^{n-k}\wedge\omega^k.
\]
No connectedness hypothesis is explicitly supplied. If disconnected manifolds are included, the same constant must satisfy the equation on every component, and integration on each component imposes the corresponding componentwise compatibility condition.

The supplied dated curation attributes solvability under a cone condition to Fang, Lai, and Ma, and says that the report leaves the case where that condition fails open. The cone condition itself is not supplied, so that residual case cannot be stated precisely from this record. These are supplied literature qualifications, not independently verified conclusions.
\subsection{Short English statement}
The question asks whether a given Kähler class contains a Kähler metric satisfying the displayed inverse \(\sigma_k\) equation. The intended range of \(k\) and the cone condition mentioned in the supplied partial-solvability assessment require clarification. This is a statement-only draft requiring human mathematical review.
\subsection{Sources}
\begin{itemize}
\item Local UnsolvedMath registry: permanent ID \(30002387\), source problem code \texttt{catalogue entry}, title \emph{Existence of Solutions to the Inverse \(\sigma_k\) Equation}; collection: Oberwolfach Reports Open Problems. No proposer is identified in the supplied record; the proposed year is recorded as 2013.
\item Supplied source citation: \emph{Partial Differential Equations (2014)}, Oberwolfach Reports, pp.~2259--2319, DOI \texttt{10.4171/owr/2013/40}. Source URL and supplied literature-source link labeled ``original OWR report'': \url{https://doi.org/10.4171/owr/2013/40}. The supplied citation uses 2014, whereas the DOI and registry year refer to 2013; this discrepancy has not been resolved.
\item Literature assessment checked at \texttt{2026-08-21}: the supplied curation labels the problem \texttt{partially\_solved}, attributes solvability under the cone condition to Fang, Lai, and Ma, and describes failure of that condition as left open by the report. This is supplied dated curation, not an independently verified literature assessment. No separate bibliographic reference for the attributed result is supplied.
\item The registry labels the statement \texttt{corrected\_verified} and reports that its mathematical content was checked against the official report, with notation, context, and question form normalized. This is a supplied verification claim; no independent source consultation or mathematical verification is asserted here.
\end{itemize}
\end{document}
